\documentclass{article}
\usepackage{graphicx}
\usepackage{amsmath}
\usepackage{amsfonts}
\usepackage{bookmark}
\usepackage{booktabs}
\usepackage{exercises}
\graphicspath{ {} }
\usepackage{indentfirst}
\usepackage{amsthm}
\usepackage{yhmath}
\usepackage{geometry}
\usepackage{placeins}
\usepackage{csquotes}
\usepackage{pgfplots}
\newcommand*{\oldneg}{\mathord{\sim}}
\newcommand\showdiv[1]{\overline{\smash{)}#1}}
\newcommand\Mydiv[1]{%
\smash{\raise.3ex\hbox{$\big)$}}\mkern-8mu
        \overline{\enspace\strut#1}}
\setlength{\tabcolsep}{5pt}
\theoremstyle{definition}
\newtheorem{definition}{Definition}

\setlength{\parindent}{.5in}
\setlength{\parskip}{12pt}
\renewcommand{\arraystretch}{2}
\geometry{margin=1in}
\newcommand{\?}{\stackrel{?}{=}}
\pgfplotsset{width=7cm,compat=1.18}\usepgfplotslibrary{statistics}

\title{6.5 Answers}
\begin{document}
\maketitle

\section*{1}
An observational study observes individuals and measures variables without controlling the individuals or their environment.

\section*{2}
Whereas an experiment occurs when treatment is put upon people in order to observe the effects of that treatment, an observational study only observes the 
effects of the treatment, and does not apply the treatment in the study.

\section*{3}
This is a biased question, as \textquote{budget cuts} on its own is a vague statement. To make the question better, specify what budget cuts were 
proposed. In addition, the question specifically mentions the cuts were proposed by the mayor. This could introduce bias as people may answer the questions just based upon whether they liked the 
mayor or not. To fix this problem, simply remove mention of the mayor.

\section*{4}
The problem here is that the option for reading a book (it says \textquote{just some book}) is made to be something lesser than the first option 
(\textquote{latest award-winning move}). This makes it more likely the first option would be chosen. To fix this problem, remove the qualifiers before listing both of the two options.

\section*{5}
The question bases itself on an unexplained premise. It states that the situation is a \textquote{health problem}, yet when describing the situation, gives no evidence 
to back up this claim. To make the question better, replace \textquote{health problem} with \textquote{situation}.

\section*{6}
The problem here is that the question is being asked from a child who is proposing the school bond issue. Thus the adult might want to please the child, and answer yes. The question could be 
improved by removing \textquote{our} from the question, and replacing it with \textquote{the}.

\section*{7}
This is an observational study. This is because the researcher is not actively making any changes to the subjects, but only observing the situation for people in 
rural and urban areas.

\section*{8}
This is an experiment, as the farmer is actively applying a new fertilizer in order to test the affects of the fertilizer.

\section*{9}
B. This is an oberservational study. It is not an experiment, as no change is being applied to the subjects. Instead, observations are being made as to the effects of 
smoking on a person's health, hence this is observational.

\section*{10-11}
10 is wrong, as the study is not just a control group. In addition, the control group described is incorrect, as it does not use neither of the conditioners, but rather the 
older one. 11 is wrong as this study involves actively applying treatment to a given group, meaning it is not observationl. The study is actually a randomized comparitive 
experiment.

\section*{12}
If we were to use experiments in astronomy, we would have to be actively applying certain conditions to the objects of our study, namely astronomical bodies. Because these 
bodies are usually things we cannot directly change, it is more likely that we would simply be observing the bodies, hence conducting observational studies.

\section*{13}
An observational study is best here, as there is not anything that could be experimented with the areas of houses. 
It could be carried out by randomly selecting some number of houses close to the park, and then randomly selecting the same number of 
houses close to the school, and then comparing the prices of the houses.

\section*{14}
This situation warrants an experiment, as treatment is being actively applied to flowers. 
It could be carried out by selecting a control group of randomly selected flowers not sprayed, and a treatment group of randomly 
selected flowers that are sprayed twice daily, and then comparing the lifespans of the flowers in the two groups (making sure the two groups are of the same size).

\section*{15}
This is not a randomized comparative experiment, as the study is conducted on 140 volunteers in each group, and hence not random people. In addition, this is not an 
experiment, but rather an observational study, as no treatment is being given to the subjects. 
Hence this study is not valid. For example, it could be that the volunteers for the group waking up after 6:30 just happened to be more dangerous drivers. 

\section*{16}
The basis of the experiment would make sense if we pick a particular breed of dogs. Then we can select a random amount of dogs from that breed. Two groups of equal size will be 
created, one given the fatty acids, the other not. Our experiment should ensure that the two groups 
live in similar conditions (i.e. amount of exercise, water, food, sleep), except for the treatment.

\section*{17}
(a): We only know that the researcher \textquote{identified} 175 people for each group. Thus, we have no guarantee that this identification process was randomized, meaning 
the results could be invalid.


(b): The study could be improved by ensuring that the subjects are chosen at random.

\section*{18}
This depends upon whether the first

\end{document}
